% Root synthesis results, complete proofs; no external file dependencies.
\section{Additional decision certificates and finite-pool checks}

\begin{proposition}[PROVED: exact zero-value and limiting conditions]
\label{prop:zero-value-r3}
If Bayes actions are attained under a fixed loss, coherent Bayes value
is zero exactly when some current Bayes action is conditionally optimal
after almost every observation. For a sequence of binary experiments,
$P(O_j=1)\to1/2$ is equivalent to $\mathbb E O_j\to0$.
For fixed finite Hamming targets with positive weights, query value tends
to zero if and only if
$(|\mathbb E[T_iO_j]|-|\mathbb E T_i|)_+\to0$ for every target.
For integrated squared graph displacement with a binary observation
whose mean tends to zero, value tends to zero if and only if
$\int w(u)\operatorname{Cov}(g(u),O_j)^2du\to0$.
\end{proposition}
\begin{proof}
For a current Bayes action $a_0$, the value equals the expectation of
the nonnegative conditional optimality gap
$\mathbb E[L(a_0,T)\mid O]-\inf_a\mathbb E[L(a,T)\mid O]$.
It vanishes exactly when that gap vanishes almost surely.
The observation-marginal equivalence follows from
$P(O=1)=(1+\mathbb E O)/2$. The Hamming result follows from the
binary-observation identity and a finite sum of nonnegative terms
with fixed positive weights. The squared-graph formula is
Proposition~\ref{rthreegeo:graphloss}; its denominator
$\operatorname{Var}(O)=1-(\mathbb E O)^2$ tends to one.
\end{proof}

\begin{proposition}[PROVED: a self-information condition for observed-label margin]
\label{prop:self-margin-r3}
Let fixed binary targets have nonnegative weights summing to one. Put
$b_j=\min\{P(O_j=1),P(O_j=-1)\}$. Then their total Hamming Bayes
value $V(j)$ is at most $b_j$. Suppose every candidate $j$ is also a target
with weight $w_j>0$, and its unweighted self-value $v_{jj}$ is at least
$\alpha b_j$ for a common $\alpha>0$. If $m$ maximizes $b_j$, then
\[
 V(m)\ge \alpha w_m\max_j V(j)\ge\alpha w_{\min}\max_j V(j).
\]
It suffices that the self-value condition hold at the selected $m$.
\end{proposition}
\begin{proof}
Let $o_+$ be a most probable observation outcome and $d_+,d_-$ the
conditional Bayes actions. Before querying, using $d_+$ is feasible,
so the old Bayes risk is at most its unconditional risk. Subtracting the
optimal observation-dependent risk leaves at most
\[
 P(O_j=o_-)\,
 \mathbb E[L(d_+,T)-L(d_-,T)\mid O_j=o_-]\le b_j,
\]
because normalized Hamming loss lies in $[0,1]$.
The selected query's total value includes its nonnegative self contribution,
so $V(m)\ge w_m v_{mm}\ge\alpha w_m b_m$. Finally
$b_m\ge b_j\ge V(j)$ for each $j$.
\end{proof}

\begin{counterexample}[COUNTEREXAMPLE: noiseless disjoint targets give no margin ratio]
Let the sole target $T$ have $P(T=1)=1/2+\delta$, where $0<\delta<1/2$.
Query $a$ reveals an independent fair sign; query $b$ reveals $T$.
Observed-label margin uniquely selects $a$, whose value is zero.
Query $b$ has value $1/2-\delta$.
\end{counterexample}
\begin{proof}
The fair observation is strictly more uncertain than the biased one.
Independence leaves the target posterior unchanged after $a$; revelation
after $b$ removes all target error.
\end{proof}

\begin{proposition}[PROVED: model-internal improvement and updater error]
\label{prop:improvement-r3}
In Theorem~\ref{thm:terminal-transfer-r3}'s notation, let
$\gamma=V_Q(j_Q)-V_Q(j_m)$, and use the $Q$-Bayes terminal rules
$g_{j_Q},g_{j_m}$. The true-risk improvement from replacing $j_m$ by
$j_Q$ is at least
\[
 \gamma-\frac{D_{j_Q}+D_{j_m}}2.
\]
If actual terminal rules $\widetilde g_j$ differ from $g_j$, define
\[
 \zeta_j=\sum_iw_iP\{\widetilde g_{ij}(O_j)\ne g_{ij}(O_j)\}.
\]
The same improvement lower bound for the actual rules subtracts the
additional term $\zeta_{j_Q}+\zeta_{j_m}$.
\end{proposition}
\begin{proof}
The fixed-rule risk discrepancy is at most $D_j/2$ by
Theorem~\ref{thm:terminal-transfer-r3}. Insert the two corresponding
$Q$ risks between the true risks. Their difference is $\gamma$ because
the baseline cancels. A classification loss can change only when the
two predicted labels differ; the absolute risk difference caused by
replacing $g_j$ is at most $\zeta_j$. Apply the triangle inequality.
\end{proof}

\begin{proposition}[PROVED: covariance and joint-log-score versions]
\label{prop:cov-kl-r3}
Write $n_j^R=\mathbb E_R O_j$ and
$k_{ij}^R=\operatorname{Cov}_R(T_i,O_j)$. If both target and observation
spin marginals differ by at most $\epsilon_m$, and
$\sum_iw_i|k_{ij}^P-k_{ij}^Q|\le\epsilon_c$ for every $j$, then
the terminal regret in Theorem~\ref{thm:terminal-transfer-r3} is at most
$2\epsilon_m+\epsilon_c$.
Alternatively, if $P_{ij},Q_{ij}$ are the actual joint laws of
$(T_i,O_j)$, then
\[
 D_j\le\sum_iw_i\sqrt{2\operatorname{KL}(P_{ij}\Vert Q_{ij})}
 \le\sqrt{2\sum_iw_i\operatorname{KL}(P_{ij}\Vert Q_{ij})}.
\]
Thus a uniform bound $K$ on the last weighted KL sum gives terminal
regret at most $\sqrt{2K}$.
\end{proposition}
\begin{proof}
Since $c_{ij}=m_i n_j+k_{ij}$ and spin means are bounded by one,
$|c_{ij}^P-c_{ij}^Q|\le |m_i^P-m_i^Q|+|n_j^P-n_j^Q|
+|k_{ij}^P-k_{ij}^Q|$. The same bound dominates the target-mean
error, proving $D_j\le2\epsilon_m+\epsilon_c$.
Each of the two moments defining $D_j$ is an expectation of a
function in $[-1,1]$, so its difference is at most $2\operatorname{TV}
(P_{ij},Q_{ij})\le\sqrt{2\operatorname{KL}(P_{ij}\Vert Q_{ij})}$
by Pinsker. Jensen's inequality for the square root gives the weighted
bound. Expected excess negative log score equals KL by its definition;
raw negative log score also contains the true entropy and is not KL.
\end{proof}

\begin{proposition}[PROVED: vanishing information under a fixed smooth channel]
\label{prop:bounded-scale-r3}
Let $G$ be a random state/field and let a query return a binary observation
with $P_c(O=1\mid G)=\kappa(cG_j)$, where $\kappa(0)=1/2$.
For any fixed measurable loss in $[0,1]$, its coherent Bayes information
value satisfies
\[
 0\le V_c(j)\le\mathbb E|\kappa(cG_j)-1/2|.
\]
Consequently the value vanishes as $c\downarrow0$ for a continuous
bounded channel. If $\mathbb E|G_j|<\infty$, logistic gives the upper
bound $c\mathbb E|G_j|/4$, and $\kappa(f)=\Phi(f/\tau)$ gives
$c\mathbb E|G_j|/(\tau\sqrt{2\pi})$.
\end{proposition}
\begin{proof}
Compare the joint law of $(G,O)$ with the law sharing the same $G$
but making $O$ an independent fair coin. Their total variation is
$\mathbb E|\kappa(cG_j)-1/2|$, as follows by summing the two outcome
probability differences conditionally on $G$. The risk of any fixed
observation-to-action rule changes by at most this TV, and taking the
infimum preserves the same bound. Under the independent coin no rule
improves on the prequery Bayes action, so its optimal risk is the
prequery risk. Nonnegativity follows from retaining that action.
Dominated convergence and the respective derivative bounds on the
links finish the proof.
\end{proof}

\begin{proposition}[PROVED: independent copy blocks make EER optimal for every budget]
\label{prop:blocks-r3}
Partition a finite target pool into blocks. Labels are identical within
each block and independent between blocks. A noiseless query reveals
the block's common sign. Let $W_b$ be total block weight and $u_b$
its initial Bayes error. Optimal total gain at budget $B$ is the sum
of the largest $\min(B,\text{number of blocks})$ numbers $W_bu_b$.
Exact greedy EER achieves it.
\end{proposition}
\begin{proof}
A query resolves its block and leaves every other block law unchanged.
Hence the value of the first query in block $b$ is always $W_bu_b$,
and additional queries there have zero value. Every adaptive
realization can resolve at most $B$ blocks, so its gain is bounded
by the sum of the largest fixed block values. Choosing those blocks
attains the bound, independently of the observed signs.
\end{proof}

\begin{counterexample}[COUNTEREXAMPLE: correct-specification multistep EER exceeds margin]
Take two independent fair singleton blocks and an independent
$M$-point copy block of error $1/2-\eta$, with equal point weights
$1/(M+2)$ and $0<\eta<1/2$. If $M(1/2-\eta)>1/2$, then at budget two
margin gains $1/(M+2)$, whereas exact EER gains
$[M(1/2-\eta)+1/2]/(M+2)$ and is horizon-optimal.
\end{counterexample}
\begin{proof}
The two fair singleton marginals are strictly most uncertain until
both are queried, so margin spends both queries there. Each contributes
$1/[2(M+2)]$. The copy-block value is strictly larger than either
singleton value; the preceding proposition gives the EER ordering
and optimum.
\end{proof}

\begin{proposition}[PROVED: randomization cannot improve a known finite-policy optimum]
For a fixed known finite-world law, finite actions/outcomes, a fixed
horizon and terminal loss, an optimal adaptive policy can be chosen
deterministic.
\end{proposition}
\begin{proof}
Condition on all internal random choices of a randomized decision tree.
Its expected risk is a mixture of risks of deterministic trees, so at
least one tree has risk no greater than the mixture. Equivalently,
backward induction minimizes an affine function of each action
distribution at an extreme point. This does not address minimax
robustness over an unknown law, where mixtures can help.
\end{proof}

\begin{proposition}[PROVED: edge-loss Bayes actions and the third moment]
\label{prop:edge-moment-r3}
On a fixed weighted graph, let $Z_e=T_iT_k$ denote whether endpoint
signs agree. For unconstrained edge-sign predictions $a_e\in\{-1,1\}$,
weighted edge-disagreement Bayes risk is
\[
 R_E=\frac12\sum_e w_e(1-|\mathbb E Z_e|).
\]
The exact one-step reduction after binary $O_j$ is
\[
 \frac12\sum_e w_e
 (|\mathbb E[T_iT_kO_j]|-|\mathbb E[T_iT_k]|)_+.
\]
If the output must be a cut induced by vertex predictions, the
optimization is over the feasible cut signs instead. Pair and triple
moments remain sufficient for evaluating and optimizing this fixed
loss before and after one binary observation, but coordinatewise edge
modes need not form a feasible cut.
\end{proposition}
\begin{proof}
Apply Lemma~\ref{lem:binary-observation-r3} to the binary target $Z_e$.
For constrained actions, the expected loss of any vertex labeling is
an affine combination of the edge means, so minimizing over finitely
many labelings uses only those means. Conditional edge means are
\[
 \mathbb E[Z_e\mid O_j=o]
 =\frac{\mathbb E Z_e+o\mathbb E[Z_eO_j]}
        {1+o\mathbb E O_j}
\]
on positive-probability outcomes. Thus second and third moments plus
the observation marginal suffice. Cycles impose compatibility
constraints: an odd number of cut edges on a cycle is impossible,
illustrating why independent edge modes may fail feasibility.
\end{proof}

\begin{counterexample}[COUNTEREXAMPLE: even an optimized ratio can have negative coherent value]
\label{ce:optimized-ratio-r3}
Use two disjoint graph edges with cut indicators $(1,Z)$, where
$P(Z=1)=p$. All four cut predictions are feasible. Consider the
posterior functional
\[
 R(p)=\min_{a\in\{0,1\}^2}
 \frac{\mathbb E\sum_{e=1}^2|Z_e-a_e|}
      {\mathbb E\sum_{e=1}^2 Z_e+\sum_{e=1}^2a_e}.
\]
At $p=3/4$, $R(p)=1/15$. A coherent binary observation with equally
probable posterior probabilities $p'=1/2$ and $p'=1$ gives expected
future risk $1/14$, hence reduction $-1/210$.
\end{counterexample}
\begin{proof}
For predictions $(0,0),(1,0),(0,1),(1,1)$ the respective ratios are
\[
 1,\qquad \frac p{2+p},\qquad\frac{2-p}{2+p},
 \qquad\frac{1-p}{3+p}.
\]
For $p\in[1/2,1]$, the last is minimal. At $3/4$ it is $1/15$,
at $1/2$ it is $1/7$, and at $1$ it is zero.
Define the observation by the joint probabilities
$P(O=0,Z=0)=1/4$, $P(O=0,Z=1)=1/4$,
$P(O=1,Z=1)=1/2$. Its two branches have the stated masses and
posteriors; no incoherent update is used. The expected future ratio
is $1/14$. Unlike this ratio of expectations, the posterior expectation
of a fixed normalized loss, minimized over fixed admissible actions,
obeys Proposition~\ref{prop:nonnegative-r3}.
\end{proof}

\begin{proposition}[PROVED: uniform first-hit and both-class discovery laws]
\label{prop:tail-r3}
For $0<r<N$, a uniform random ordering of $N$ items with $r$ designated
rare items has first-hit time $D$ satisfying, for $0\le k\le N$,
\[
 P(D>k)=\frac{\binom{N-r}{k}}{\binom Nk},\qquad
 \mathbb E D=\frac{N+1}{r+1}.
\]
Its time $T$ to observe both classes satisfies, for $1\le k\le N$,
\[
 P(T>k)=\frac{\binom rk+\binom{N-r}k}{\binom Nk},\qquad
 \mathbb E T=\frac{N+1}{r+1}+\frac{N+1}{N-r+1}-1.
\]
Binomial coefficients exceeding available items are zero; survivors
are zero beyond the population size.
In a fixed subset with $M$ items and $s\ge1$ rare items, uniform
ordering has the same first-hit formulas with $N,r$ replaced by $M,s$.
If $s/M\ge r/N$, this subset first-hit time stochastically dominates
in the favorable direction: its survival probability is no larger.
\end{proposition}
\begin{proof}
No rare item among the first $k$ is precisely a size-$k$ subset chosen
from the $N-r$ common items, giving the survival formula.
The $N-r$ common items occupy $r+1$ exchangeable gaps around the rare
items, so the expected number before the first rare is $(N-r)/(r+1)$.
For $T>k$ with $k\ge1$, all $k$ labels are rare or all are common,
disjoint events. Summing their survivor probabilities gives the
displayed mean, with the duplicated $k=0$ term corrected by subtracting
one. Apply the same argument to the subset.
Conditional on $k$ failures, the subset success hazard is $s/(M-k)$
and the full-pool hazard is $r/(N-k)$ while both are defined.
The difference after cross multiplication has numerator
$sN-rM+k(r-s)\ge0$, because $s\le r$ and enrichment holds.
Multiplying failure hazards proves the survivor ordering; after
the subset's common items are exhausted its survival is zero.
Neither argument applies to an arbitrary history-conditioned
deterministic extreme-selection policy.
\end{proof}

\begin{counterexample}[COUNTEREXAMPLE: no arbitrary continuous-boundary guarantee from finite signs]
\label{ce:finite-boundary-r3}
Let a finite query pool in an interval leave an open gap $(a,b)$.
Two noiseless smooth classifiers $f_0(x)=x-\theta_0$ and
$f_1(x)=x-\theta_1$, with $a<\theta_0<\theta_1<b$, agree on every
pool label but have boundary Hausdorff distance $\theta_1-\theta_0$.
Any pool-only estimator has worst-world expected boundary error at
least $(\theta_1-\theta_0)/2$.
\end{counterexample}
\begin{proof}
Every pool point is on the same side of both thresholds. Thus even
revealing all pool labels gives identical data in the two worlds.
Couple the estimator's randomness so its output is the same.
The triangle inequality for Hausdorff distance implies that its two
errors sum to at least the distance between the boundaries. Take
expectations and the larger error. The thresholds can approach the
gap endpoints. A coverage bound limits this ambiguity but an arbitrary
finite pool supplies no smaller universal positional guarantee.
\end{proof}
